% Partie : méthodes
% Chapitre : induction mathématique
% Section : introduction

\documentclass[../../../include/open-logic-section]{subfiles}

\begin{document}
\olfileid[fr]{mth}{ind}{int}

\olsection{Introduction}

L'induction mathématique est une
méthode de démonstration importante
qui, sous différentes formes,
s'emploie dans presque tous les
domaines de la logique, de
l'informatique théorique et des
mathématiques. Elle est nécessaire
pour démontrer de nombreux
résultats de logique.

On oppose souvent l'induction
à la déduction et on la décrit
comme le passage du particulier
au général. Supposons, par
exemple, que nous ayons observé
de nombreuses émeraudes vertes
et aucune émeraude qui ne soit
pas verte. Nous pourrions en
conclure que toutes les
émeraudes sont vertes.
Cette inférence est inductive :
elle passe de nombreux cas
particuliers (cette émeraude
est verte, celle-là est verte,
etc.) à un énoncé général
(toutes les émeraudes sont
vertes). L'induction
\emph{mathématique} conclut,
elle aussi, à un énoncé
général, mais elle est
très différente de cette
«~induction simple~».

En termes très généraux,
une démonstration par
induction mathématique
établit que tous les
objets mathématiques
d'un certain type
possèdent une propriété.
Dans le cas le plus
simple, ces objets
sont les nombres
naturels. On établit
alors qu'ils possèdent
tous une propriété en
montrant que :
\begin{enumerate}
    \item $0$ possède la propriété ;
    \item chaque fois qu'un nombre~$k$
      la possède, $k+1$ la possède aussi.
\end{enumerate}
Le raisonnement par récurrence
sur les nombres naturels
permet souvent de démontrer
des énoncés généraux
concernant d'autres objets
mathématiques auxquels
on peut associer un nombre.
Par exemple, chaque ensemble
fini possède un nombre fini~$n$
d'éléments. Si l'on démontre
par récurrence que, pour
chaque~$n$, «~tous les
ensembles finis de cardinal~$n$
sont $\dots$~», on obtient
un résultat sur tous
les ensembles finis.

L'induction s'étend aussi
aux objets mathématiques
qui sont \emph{définis
inductivement}. Ainsi,
les expressions d'un
langage formel, comme
celles de la logique
du premier ordre,
reçoivent une définition
inductive. L'\emph{induction
structurelle} permet de
démontrer des résultats
concernant toutes ces
expressions. Elle est
particulièrement utile
et très employée
en logique.

\end{document}
